\documentclass[12pt,a4paper]{extarticle}
\usepackage{cmap}
\usepackage[T2A]{fontenc}
\usepackage[utf8]{inputenc}
\usepackage[english,russian]{babel}
\usepackage[left=30mm,right=15mm,top=20mm,bottom=20mm]{geometry}
\usepackage{setspace}
\onehalfspacing
\usepackage{indentfirst}
\setlength{\parindent}{1.25cm}
\usepackage{longtable}
\usepackage{array}
\usepackage{listings}
\usepackage{xcolor}
\usepackage{amsmath}
\usepackage{titlesec}
\usepackage{caption}
\captionsetup[longtable]{labelsep=endash,justification=raggedright,
  singlelinecheck=false,font=normalsize}
\usepackage{enumitem}
\usepackage{xurl}
\usepackage[hidelinks,unicode]{hyperref}

\urlstyle{tt}
\newcommand{\code}[1]{\nolinkurl{#1}}
\newcommand{\struct}[1]{%
  \clearpage\begingroup
  \titlespacing*{\section}{0pt}{0pt}{18pt}%
  \section*{\centering #1}\endgroup%
  \addcontentsline{toc}{section}{#1}}
\newcommand{\mainsection}[1]{\clearpage\section{#1}}
\titleformat{\section}{\normalfont\normalsize}{\thesection}{1em}{}
\titleformat{\subsection}{\normalfont\normalsize}{\thesubsection}{1em}{}
\titleformat{\subsubsection}{\normalfont\normalsize}{\thesubsubsection}{1em}{}
\titlespacing*{\section}{\parindent}{0pt}{18pt}
\titlespacing*{\subsection}{\parindent}{18pt}{12pt}
\titlespacing*{\subsubsection}{\parindent}{12pt}{6pt}
\setcounter{tocdepth}{2}
\setlist{nosep,leftmargin=1.25cm}
\setlength{\tabcolsep}{3pt}
\renewcommand{\arraystretch}{1.15}
\setlength{\LTleft}{0pt}
\setlength{\LTright}{\fill}
\setlength{\LTcapwidth}{\textwidth}
\newcolumntype{P}[1]{>{\raggedright\arraybackslash}p{#1}}
\renewcommand{\contentsname}{СОДЕРЖАНИЕ}
\renewcommand{\tablename}{Таблица}
\AtBeginDocument{\renewcommand{\thetable}{\arabic{table}}}
\lstdefinelanguage{Go}{%
  morekeywords={package,import,func,type,struct,interface,return,if,else,for,range,var,const,nil,true,false,string,int,int64,bool},%
  sensitive=true,morecomment=[l]{//},morestring=[b]"}
\lstset{language=Go,basicstyle=\ttfamily\small,
  keywordstyle=\color{black},commentstyle=\color{black!65},
  stringstyle=\color{black},extendedchars=true,inputencoding=utf8,
  columns=fullflexible,keepspaces=true,showstringspaces=false,
  breaklines=true,breakatwhitespace=false,tabsize=4,
  frame=single,rulecolor=\color{black},numbers=none}
\emergencystretch=3em
\widowpenalty=10000
\clubpenalty=10000
\pagestyle{plain}

% Six attributes: the first appears in the heading, five in the table.
\newcommand{\testcase}[6]{%
  \par\addvspace{12pt}\noindent #1\par\nopagebreak
  \begin{longtable}{|P{0.22\textwidth}|P{0.72\textwidth}|}
  \hline
  Тип теста & #2 \\ \hline
  Объект & #3 \\ \hline
  Входные данные & #4 \\ \hline
  Косвенные входные данные & #5 \\ \hline
  Ожидаемый результат и проверка & #6 \\ \hline
  \end{longtable}}

\begin{document}
\hypersetup{pageanchor=false}
\begin{titlepage}
\centering
Учебный проект по тестированию программного обеспечения\par
\vfill
ПЛАН ТЕСТИРОВАНИЯ\par
веб-приложения «Задачи» на языке Go\par
\vspace{1cm}
Версия плана: 1.0\par
Состояние: до реализации приложения и выполнения тестов\par
\vfill
\begin{flushright}
Исполнитель: \underline{\hspace{5cm}}\\[6pt]
Группа: \underline{\hspace{5cm}}\\[6pt]
Проверяющий: \underline{\hspace{5cm}}
\end{flushright}
\vfill
\underline{\hspace{4cm}}\par
2026
\end{titlepage}
\setcounter{page}{2}
\hypersetup{pageanchor=true}

\struct{РЕФЕРАТ}
Объект планируемого исследования --- однопользовательское веб-приложение
«Задачи». Цель документа --- задать воспроизводимые проверки до написания
приложения и определить критерии его приемки. Устанавливаются контракты
HTTP API, бизнес-правила, схема хранения и идентификаторы элементов интерфейса.
Предусмотрены блочные, интеграционные, аттестационные, нагрузочные проверки,
а также проверки SQL-инъекций и XSS.

Документ содержит шесть нумерованных разделов и 20 спецификаций тестов:
по пять на каждый уровень. Параметризованные спецификации содержат несколько
обязательных наборов данных. Журналы и показатели разделов 4--6 оставлены
для заполнения после проведения тестирования. Числа в критериях являются
плановыми порогами, а не результатами измерений.

Ключевые слова: Go, SQLite, HTTP API, сервис задач, тестирование,
покрытие операторов, покрытие ветвей, аттестация.

\clearpage
\begingroup
\titleformat{\section}{\normalfont\normalsize\centering}{}{0pt}{}
\titlespacing*{\section}{0pt}{0pt}{18pt}
\tableofcontents
\endgroup

\struct{ВВЕДЕНИЕ}
Приложение предназначено для ведения небольшого личного списка задач
на локальном компьютере. Пользователь создает задачу, просматривает список,
отмечает задачу выполненной и удаляет запись. Один сервер обслуживает один
общий список; учетных записей и разграничения доступа нет.

План является основанием для последующей реализации. Изменение интерфейса,
бизнес-правила или кода ответа требует изменения версии плана и связанных
проверок. Предполагаемый объем прикладного кода Go --- 300--500 непустых строк
без тестов, зависимостей, HTML и JavaScript. Это ограничение проекта,
а не измеренный объем существующей программы.

Оформление ориентировано на ГОСТ 7.32-2001 с изменением №~1;
параметры задания: A4, 12 пт, полуторный интервал, поля 30/15/20/20 мм,
абзац 1,25 см. Используется именно указанная в задании редакция стандарта.
Документ собирается командой \code{pdflatex task_test_plan.tex} дважды
для формирования содержания; внешние файлы, BibTeX и shell-escape не нужны.

\mainsection{Объект тестирования}
\subsection{Назначение и состав приложения}
«Задачи» --- серверное приложение с HTML-страницей по адресу
\url{http://127.0.0.1:8080/}. Сервер использует Go, \code{net/http},
\code{encoding/json}, \code{html/template}, \code{database/sql}.
Хранилище --- файловая SQLite с драйвером
\code{github.com/mattn/go-sqlite3}; сборка требует CGO и компилятора C.
Страница содержит форму создания и таблицу задач; JavaScript вызывает API
через \code{fetch}. Начальные записи выводятся шаблоном, последующие изменения
DOM --- через \code{textContent}, без вставки названия через \code{innerHTML}.

Сервер слушает только \code{127.0.0.1:8080}. Сетевой доступ других компьютеров,
публикация в Интернете и режим нескольких пользователей не входят в проект.
Файл БД задается параметром \code{DB_PATH}; при отсутствии файла создается
схема, при следующем запуске существующие записи сохраняются.

\subsection{Данные и бизнес-правила}
Модель \code{model.Task} имеет поля \code{ID int64}, \code{Title string},
\code{Done bool}; JSON-имена: \code{id}, \code{title}, \code{done}.
Других полей в ответе задачи нет.

\begin{enumerate}
\item Название сначала обрабатывается \code{strings.TrimSpace}.
Затем проверяются UTF-8 и длина \code{utf8.RuneCountInString}:
$1\leq L\leq100$. Единица длины --- кодовая точка Unicode,
а не байт и не визуальный графемный кластер. Внутренние пробелы сохраняются.
Одинаковые названия разрешены. HTML и SQL в названии хранятся как текст.
\item ID --- целое $I\in[1;2147483647]$. В URL используется десятичная
запись без знака, пробелов и ведущих нулей. Значения \code{0}, \code{-1},
\code{01}, \code{abc}, \code{2147483648} отклоняются до вызова репозитория.
Автоматический ID БД не переиспользуется после удаления.
\item Новая задача получает \code{Done=false}.
Отметка выполнения устанавливает \code{Done=true}; повторная отметка
возвращает ту же задачу с \code{Done=true}. Обратного переключения нет.
\item Список отсортирован по \code{id ASC}; пустой список --- JSON-массив
\code{[]}, а не \code{null}. Удаление существующей записи физическое;
повторное удаление возвращает признак отсутствия записи.
\item При отказе проверки ввода репозиторий не вызывается.
Сбой хранилища не раскрывает клиенту путь файла, SQL или текст драйвера.
\end{enumerate}

Схема и запросы будущего репозитория:
\begin{lstlisting}[language=SQL]
CREATE TABLE IF NOT EXISTS tasks (
  id INTEGER PRIMARY KEY AUTOINCREMENT
     CHECK (id BETWEEN 1 AND 2147483647),
  title TEXT NOT NULL CHECK (length(title) BETWEEN 1 AND 100),
  done INTEGER NOT NULL DEFAULT 0 CHECK (done IN (0, 1))
);
-- Values are passed as database/sql arguments.
INSERT INTO tasks(title, done) VALUES (?, 0);
SELECT id, title, done FROM tasks ORDER BY id ASC;
UPDATE tasks SET done = 1 WHERE id = ?;
DELETE FROM tasks WHERE id = ?;
\end{lstlisting}

Проверка Unicode выполняется сервисом; ограничение SQLite служит
дополнительным барьером и не заменяет правила сервиса.
Для SQLite устанавливаются \code{SetMaxOpenConns(1)},
\code{PRAGMA busy_timeout=5000}, \code{PRAGMA journal_mode=WAL}.
Метод завершения выполняет UPDATE и SELECT в одной транзакции,
возвращает задачу после фиксации. Нулевое число затронутых строк у UPDATE
или DELETE преобразуется в \code{ErrNotFound}.

\subsection{Контракт HTTP и пользовательского интерфейса}
В таблице~\ref{tab:api} задан транспортный контракт.
У POST заголовок \code{Content-Type: application/json}; параметр
\code{charset=utf-8} допускается. JSON содержит только строковое поле
\code{title}; отсутствующее поле, \code{null}, иной тип, лишнее поле,
поврежденный JSON и второе JSON-значение дают \code{400 bad_json}.
Декодер использует запрет неизвестных полей и проверку EOF после первого
значения. PATCH и DELETE отправляются без тела.

\begin{longtable}{|P{43mm}|P{30mm}|P{79mm}|}
\caption{Контракт API}\label{tab:api}\\\hline
Метод и путь & Код & Тело и действие \\\hline\endfirsthead
\multicolumn{3}{l}{Продолжение таблицы \thetable}\\\hline
Метод и путь & Код & Тело и действие \\\hline\endhead
POST \code{/api/tasks} & 201 & Объект задачи; например,
\code{{"id":1,"title":"Отчет","done":false}} \\\hline
GET \code{/api/tasks} & 200 & Массив объектов задач по \code{id ASC} \\\hline
PATCH \code{/api/tasks/1/done} & 200 &
\code{{"id":1,"title":"Отчет","done":true}} \\\hline
DELETE \code{/api/tasks/1} & 204 & Тело длиной 0 байт \\\hline
\end{longtable}

Все ответы с JSON имеют \code{Content-Type: application/json; charset=utf-8}.
Сравнение JSON выполняется после декодирования: порядок ключей,
пробелы и завершающий перевод строки не учитываются; набор ключей,
типы и значения должны совпасть. Для DELETE 204 заголовок Content-Type
не проверяется. Страница GET / возвращает 200 и
\code{text/html; charset=utf-8}.

Тело отказа имеет ровно одно поле: \code{{"error":"код"}}.
Соответствия: \code{ErrTitle} --- 400 и \code{title_length};
\code{ErrID} --- 400 и \code{bad_id}; \code{ErrNotFound} --- 404 и
\code{not_found}; \code{ErrStorage} --- 500 и \code{storage}.
Неподдерживаемый Content-Type у POST дает 415 и \code{media_type}.
Метод PUT на \code{/api/tasks} дает 405 и \code{method},
\code{Allow: GET, POST}. Неизвестный путь \code{/api/unknown}
дает 404 и \code{route}. Для ошибок запроса действует порядок:
маршрут/метод, Content-Type, JSON, бизнес-правила, хранилище.

На странице заданы \code{\#title-input}, \code{\#create-button},
\code{\#error-banner}, \code{\#tasks-body}. Строка задачи имеет
\code{data-id="1"}; ячейки --- классы \code{title}, \code{status};
кнопки --- \code{.done-button}, \code{.delete-button}.
Текст статуса: «Не выполнена» или «Выполнена».
Кнопка завершения блокируется после \code{Done=true}.
Перед запросом баннер очищается. При ответе 400 с \code{title_length}
баннер содержит ровно «Название: от 1 до 100 символов»;
при 500 с \code{storage} --- «Хранилище недоступно».
После создания поле ввода очищается, строка появляется в таблице.
После удаления строка отсутствует. Обновление DOM завершается не позднее
2 с после получения ответа; в Playwright применяются ожидания состояния,
а не фиксированные паузы.

\subsection{Рамки и функциональные требования}
Включаются исходный прикладной код, HTML/JS, HTTP-контракт, SQL-запросы,
отображение ошибок, сохранение после перезапуска и защита точек ввода.
Не проверяются внутренние алгоритмы стандартной библиотеки Go,
драйвера, SQLite и браузера; их совместная работа с приложением проверяется.
Исключаются аутентификация, роли, TLS, облачное развертывание, резервное
копирование, мобильная верстка, доступность и произвольные уровни нагрузки.

\begin{longtable}{|P{12mm}|P{30mm}|P{65mm}|P{43mm}|}
\caption{Функциональности и участие в аттестации}\\\hline
ID & Название & Проверяемое требование & Аттестация и обоснование \\\hline\endfirsthead
\multicolumn{4}{l}{Продолжение таблицы \thetable}\\\hline
ID & Название & Проверяемое требование & Аттестация и обоснование \\\hline\endhead
F01 & Создание & Название сохраняется после TrimSpace; ID в заданном
диапазоне, Done=false & Да: A01, основной ввод пользователя \\\hline
F02 & Просмотр & Все записи по ID ASC, пустой список [] & Да: A02,
наблюдаемый список \\\hline
F03 & Завершение & Done=true; повтор не меняет запись & Да: A03,
изменение состояния \\\hline
F04 & Удаление & Строка исчезает из БД и страницы & Да: A04,
завершение жизненного цикла \\\hline
F05 & Проверки и отказы & Ограничения названия/ID; коды HTTP и баннеры
по контракту & Да: A01, A04, A05, наблюдаемые отказы \\\hline
F06 & Долговременное хранение & Задача и Done сохраняются после
остановки процесса и открытия того же файла & Да: A03,
пользователь продолжает работу \\\hline
\end{longtable}

Нефункциональные требования: N01 --- нагрузочные пороги S01--S02;
N02 --- SQL остается данными в S03; N03 --- название не исполняется
как HTML/JavaScript в S04--S05. Их проверяют специальные тесты.

\mainsection{Стратегия тестирования}
\subsection{Структура, связи и контракты слоев}
\code{cmd/tasks/main.go} открывает БД, создает репозиторий, сервис,
обработчики и HTTP-сервер; содержит только сборку зависимостей.
Пакеты: \code{internal/model}, \code{internal/repository},
\code{internal/service}, \code{internal/handler}, \code{web}.
Последний содержит HTML-шаблон и JavaScript, встраиваемые в бинарный файл.

Направление вызовов при изменении: браузер вызывает handler;
handler вызывает service; service вызывает интерфейс Repository;
SQLiteRepository вызывает database/sql и SQLite.
При GET / обработчик получает список через тот же сервис и передает
его в html/template. \code{model} импортируется слоями, но сам
не зависит от них. Ни handler, ни web не выполняют SQL.

Публичные контракты сервиса и репозитория:
\begin{lstlisting}
// internal/model
// Task JSON tags are id, title, done.
type Task struct {
    ID int64 `json:"id"`
    Title string `json:"title"`
    Done bool `json:"done"`
}

// internal/service; imports context and internal/model.
type Repository interface {
    Create(context.Context, string) (model.Task, error)
    List(context.Context) ([]model.Task, error)
    Complete(context.Context, int64) (model.Task, error)
    Delete(context.Context, int64) error
}
// Constructor: New(repo Repository) *Service
// Methods of *Service:
// Create(ctx context.Context, title string) (model.Task, error)
// List(ctx context.Context) ([]model.Task, error)
// Complete(ctx context.Context, id int64) (model.Task, error)
// Delete(ctx context.Context, id int64) error
\end{lstlisting}

\code{ErrTitle}, \code{ErrID}, \code{ErrNotFound}, \code{ErrStorage}
объявляются в model как переменные \code{errors.New}.
При отказе методы, возвращающие Task, возвращают \code{model.Task{}}
и ошибку, распознаваемую \code{errors.Is}; List возвращает
\code{nil, err}. При отсутствии задач List возвращает срез \code{[]model.Task{}}
и \code{nil}: \code{tasks != nil}, \code{len(tasks) == 0}.

\begin{longtable}{|P{35mm}|P{60mm}|P{57mm}|}
\caption{Участие структурных элементов в проверках}\\\hline
Элемент & Уровни и объекты & Окружение или замена \\\hline\endfirsthead
\multicolumn{3}{l}{Продолжение таблицы \thetable}\\\hline
Элемент & Уровни и объекты & Окружение или замена \\\hline\endhead
model & Структуры и ошибки используются всеми тестами;
самостоятельной логики нет & Эталонные Task и errors.Is \\\hline
service & Блочные U01--U05; интеграционные I02--I05;
аттестационные A01--A05 & В U: ручная заглушка Repository \\\hline
repository & Интеграционные I01--I05; косвенно A и S & Реальная SQLite
в отдельном временном файле; драйвер database/sql \\\hline
handler & Интеграционные I02--I05; аттестационные A01--A05 &
httptest и настоящий сервис; при отказе БД --- закрытый sql.DB \\\hline
web & Аттестационные A01--A05; специальные S04--S05 &
Chromium, Playwright, настоящий сервер \\\hline
cmd/tasks & Сборка зависимостей и перезапуск A03 &
Отдельный процесс; без подмен слоев \\\hline
Go, SQLite, ОС & Окружение всех уровней & Не оцениваются по покрытию
прикладного кода \\\hline
\end{longtable}

\subsection{Блочное тестирование}
Метод --- белый ящик для развилок сервиса и черный ящик для его
публичного контракта. Классы названий: после TrimSpace длина 0,
1--100 и более 100; дополнительно строка с байтом \code{0xff}.
Границы: 0, 1, 100, 101. Классы ID: меньше 1, от 1 до верхней границы,
выше верхней границы. Развилки репозитория: nil, ErrNotFound, ErrStorage.

Инструменты: \code{testing}, табличные подтесты \code{t.Run}, ручные
заглушки с подсчетом вызовов. Применение testify и генератора gomock
для данного объема не требуется. Аналог require реализуется через
\code{t.Fatalf}: сравниваются все поля Task, длины/элементы срезов,
число вызовов, аргументы и \code{errors.Is}. Текст ошибки драйвера
не является эталоном. У каждого подтеста собственная заглушка;
реальная БД и HTTP не используются.

\subsection{Интеграционное тестирование}
Схема --- восходящая. Проверяются контракты между реальными слоями,
SQL-параметры, фиксация изменений, JSON и отображение ошибок в HTTP.
Последовательность:
\begin{enumerate}
\item Repository + database/sql + SQLite: I01.
Драйвер теста вызывает методы напрямую; сервис и handler отсутствуют,
заглушек нет.
\item Service + Repository + SQLite: I02.
Драйвер теста вызывает сервис и проверяет БД независимыми SELECT;
handler отсутствует, заглушек нет.
\item Handler + Service + Repository + SQLite: I03--I05.
Драйверы --- \code{httptest.NewRequest}, \code{NewRecorder}
и при сетевом запуске \code{httptest.NewServer}.
В I05 после фикстуры вызывается \code{db.Close()}; это управляемый отказ
реального соединения, а не замена репозитория моком.
\end{enumerate}

Каждый независимый вариант получает отдельный файл
\code{filepath.Join(t.TempDir(), "tasks.db")}. Проверки выполняются
последовательно. Фикстура F0 --- только что созданная пустая схема:
таблица пуста, таблица \code{sqlite_sequence} не содержит строки tasks;
первый INSERT получает ID=1. Фикстура F1 --- F0 плюс SQL:
\begin{lstlisting}[language=SQL]
INSERT INTO tasks(id,title,done)
VALUES (1, CAST(X'D09ED182D187D0B5D182' AS TEXT), 0);
\end{lstlisting}
В SQL-листинге название «Отчет» задано UTF-8-байтами через hex-литерал;
в листинге Go используется Unicode-escape.
Данные других тестов не переносятся; операционный файл пользователя
не открывается. Эталоны сверяются через database/sql на том же файле
после завершения запроса, а не через методы проверяемого репозитория.

\subsection{Аттестационное тестирование}
Метод --- черный ящик: пользовательские сценарии через браузер
и контроль API независимым HTTP-клиентом Playwright.
Запускаются бинарный файл сервера, файловая SQLite и настоящий frontend;
моки HTTP и Repository запрещены. \code{curl} может применяться для
воспроизведения запросов; автоматизированный эталон --- Playwright.
Перед каждым A-тестом создается F0 или F1 в отдельном каталоге.

Playwright проверяет значения \code{textContent}, число DOM-узлов,
\code{disabled}, URL запроса, статус и JSON сетевого ответа.
Тайм-аут сетевого ожидания --- 5 с; DOM после ответа --- 2 с.
В A03 процесс останавливается с ожиданием завершения не более 5 с,
затем запускается с тем же \code{DB_PATH}. Старт подтверждается GET /api/tasks
с кодом 200 в течение 5 с. Другие A-тесты используют новый файл.

\subsection{Специальные виды тестов}
\subsubsection{Нагрузка}
Инструмент --- k6 с \code{constant-arrival-rate}, один запрос за итерацию,
без sleep и повторных запросов. Отдельные прогоны чтения и записи
описаны S01--S02. Прогрев 10 с выполняется отдельным прогоном и не
включается в метрики; затем восстанавливается начальная фикстура.
Для измерительного прогона: rate=50, timeUnit=1s, duration=60s,
preAllocatedVUs=20, maxVUs=100. Плановое число запросов --- 3000.

Пороги для каждого прогона: достигнутая скорость $N/60\geq49$ запросов/с,
$N\geq2940$, \code{dropped_iterations=0},
\code{p(95)<=200} мс для \code{http_req_duration}, число 5xx равно 0,
\code{http_req_failed=0}, доля выполненных k6 checks равна 1.
Каждый check проверяет код и JSON-схему, для записи также title и done.
Число 5xx собирается отдельным Counter; отсутствие метрики с образцами
трактуется как ноль только при подтвержденном $N\geq2940$.

Опорный стенд: Linux x86-64, серверу доступны 2 логических CPU,
2 ГиБ RAM, локальный SSD; генератор --- другой процесс на выделенных
2 CPU того же хоста, общая RAM хоста не менее 8 ГиБ. Параллельные
задачи сборки отключены. Перед прогоном фиксируются модель CPU,
ОС, версии инструментов и размер БД. Пороги относятся к этому стенду;
при смене ресурсов условия отдельно согласуются до прогона.

\subsubsection{Безопасность}
S03 проверяет параметризацию SQL: вредоносная строка хранится
побайтово, таблица tasks остается доступной. S04 проверяет
построение DOM после fetch; S05 --- HTML начальной загрузки через
html/template. Перед навигацией Playwright через init-script
задает \code{window.__xss=0}; после операций счетчик должен оставаться 0.
Фиксируются также появление диалогов и неожиданных элементов.
Проверяется отсутствие преобразования пользовательского текста в
\code{template.HTML}. Это проверки конкретных точек ввода;
они не являются заявлением об отсутствии всех классов уязвимостей.

\subsection{Стенд, артефакты и оценивание отдельного теста}
Базовый стек стенда: Go 1.23.x, \code{CGO_ENABLED=1}, SQLite через
зафиксированный в go.mod драйвер, Chromium через Playwright, k6.
Точные patch-версии Go, драйвера и встроенной SQLite, ОС, Node.js,
Playwright, Chromium и k6 записываются до первого прогона.
\code{go.sum} и lock-файл JavaScript сохраняются;
во всех сравниваемых прогонах используются те же версии.

Артефакты прогона: хеш коммита приложения, версия плана,
команды, дата/время с часовым поясом, stdout/stderr Go,
coverage.out, coverage.html, отчет Playwright с trace,
сводка k6 в JSON и файл БД теста при дефекте.
Статусы: PASS --- все сравнения данного набора совпали;
FAIL --- хотя бы одно сравнение не совпало; BLOCKED --- проверка
не выполнена из-за среды или зависимости; NOT RUN --- не запускалась.
Ни один из этих статусов заранее не присваивается тестам.

\subsection{Начало, окончание и переходы}
\begin{enumerate}
\item Начало блочного этапа: реализованы контракты подраздела 2.1,
\code{go build ./...} завершается с кодом 0; U-тесты компилируются;
зафиксированы версия плана, commit и зависимости.
\item Блочное $\rightarrow$ интеграционное: все наборы U01--U05 имеют
PASS; покрытие операторов service не ниже 90\%; обе стороны каждого
условия проверки названия/ID и обработки ошибки Repository выполнены
и отмечены в реестре ветвей; открытых блокирующих и критических дефектов 0.
\item Интеграционное $\rightarrow$ аттестационное: все наборы I01--I05
имеют PASS; общий порог покрытия подраздела~\ref{sec:coverage} достигнут;
\code{go test -race ./internal/...} дает код 0 и 0 сообщений
\code{DATA RACE}; открытых блокирующих и критических дефектов 0.
\item Аттестационное $\rightarrow$ специальное: все наборы A01--A05
имеют PASS на одном commit; F01--F06 покрыты аттестацией на 100\%;
открытых блокирующих и критических дефектов 0.
\item Окончание: все наборы U, I, A, S имеют PASS после последних
изменений, пороги S01--S02 соблюдены, S03--S05 не фиксируют исполнения
кода/SQL-команд из названий, выполнены критерии раздела 6.
\end{enumerate}

Один этап может содержать повторные запуски после исправлений.
Изменение исходного кода аннулирует относящиеся к нему результаты;
повторяются затронутые тесты и U/I-регрессия, затем A/S на итоговом commit.
Перекрытие этапов для диагностики разрешено, но не заменяет критерий перехода.

\subsection{Приостановка, возобновление и дефекты}
Выполнение затронутого этапа приостанавливается при любом событии:
сборка завершилась с ненулевым кодом; создание/открытие тестовой БД
или \code{PingContext} с тайм-аутом 2 с вернуло error;
сервер завершился вне команды теста; GET /api/tasks недоступен
три раза подряд с интервалом 1 с; обнаружено изменение файла
вне временного каталога; базовая операция создания теряет запись;
свободное место стало менее 1 ГиБ. Запланированный \code{db.Close()}
в I05 и перезапуск A03 не считаются сбоем среды.

Возобновление: причина устранена; сборка имеет код 0;
на новой F0 два последовательных контрольных цикла
POST «Контроль» $\rightarrow$ GET $\rightarrow$ DELETE возвращают
201, 200, 204, JSON содержит созданную задачу, а SELECT после DELETE
возвращает COUNT(*)=0. Полученные ID берутся из POST, поскольку
после первого цикла следующий ID уже не равен 1.
Заблокированные и затронутые тесты запускаются заново.

Блокирующий дефект: сервер не запускается или тесты не могут
продолжаться. Критический: потеря/изменение чужой записи, выполнение
внедренного SQL/JavaScript, нарушение F01--F04 без обходного пути.
Существенный: нарушение иного контракта без потери данных.
Незначительный: оформление при сохранении контрактов и возможности работы.
Классификация записывается в описании дефекта, а не выводится из кода HTTP.

\mainsection{Детальный план тестов}
\subsection{Общие данные и порядок выполнения}
Префиксы U, I, A, S обозначают блочные, интеграционные,
аттестационные и специальные проверки. Цель и ID составляют
первый атрибут каждой спецификации; остальные пять атрибутов
приведены в ее таблице. Все перечисленные варианты обязательны.
Каждый параметризованный вариант независим, если явно не задана
последовательность действий внутри одного сценария.

T100 --- строка из ровно 100 ASCII-букв A, T101 --- из 101 буквы A.
Их точные литералы (перенос при печати не является частью строки):
\begin{quote}\raggedright
T100 = \code{AAAAAAAAAAAAAAAAAAAAAAAAAAAAAAAAAAAAAAAAAAAAAAAAAAAAAAAAAAAAAAAAAAAAAAAAAAAAAAAAAAAAAAAAAAAAAAAAAAAA}\par
T101 = \code{AAAAAAAAAAAAAAAAAAAAAAAAAAAAAAAAAAAAAAAAAAAAAAAAAAAAAAAAAAAAAAAAAAAAAAAAAAAAAAAAAAAAAAAAAAAAAAAAAAAAA}
\end{quote}
Для U/I контекст --- \code{context.Background()}, если не указано иное.
У A/S используется origin \url{http://127.0.0.1:8080};
сессий и Cookie нет. У JSON-запросов POST по умолчанию задан
\code{Content-Type: application/json}; у GET --- тело отсутствует.
Для каждого теста с F0/F1 применяется новый файл БД.
Нормы $L$ и $I$ определены в подразделе 1.2 и действуют во всех тестах.

\subsection{Блочные тесты}
\testcase{U01. Цель: создание передает нормализованное название репозиторию (F01)}
{Общий}
{\code{service.(*Service).Create}}
{\code{title="  Отчет  "}; после TrimSpace L=5, диапазон 1--100.}
{Заглушка Create при аргументе «Отчет» возвращает
\code{Task{ID:1, Title:"Отчет", Done:false}, nil}; остальные методы
не вызываются. БД отсутствует.}
{Возвращены именно указанный Task и \code{err==nil}; Create вызван
один раз с побайтовой строкой «Отчет». Проверка через testing:
сравнение полей, строки аргумента и счетчика вызовов.}

\testcase{U02. Цель: проверить границы названия и UTF-8 (F01, F05)}
{Краевой; варианты с длиной вне диапазона и байтом 0xff --- негативные}
{\code{service.(*Service).Create}}
{Отдельные наборы: \code{"Я"} (L=1), T100 (L=100),
\code{""} (L=0), \code{"   "} (после TrimSpace L=0), T101 (L=101),
\code{string([]byte{0xff})} (не UTF-8). Допуск: UTF-8, L=1--100.}
{Новая заглушка для каждого набора. Для «Я» и T100 Create возвращает
Task с ID=1, тем же названием, Done=false и nil.
Для остальных наборов вызов заглушки запрещен.}
{Для первых двух: Task с указанными полями, nil, счетчик Create=1.
Для остальных: \code{Task{}}, \code{errors.Is(err,model.ErrTitle)==true},
счетчик=0. Все шесть наборов сравниваются в t.Run.}

\testcase{U03. Цель: границы ID и запрет обращения к данным при выходе из диапазона (F03--F05)}
{Краевой и негативный}
{\code{service.(*Service).Complete}, \code{Delete}}
{Каждый метод с ID=1, 2147483647, 0, -1, 2147483648;
допуск I=1--2147483647.}
{Для каждого метода/ID новая заглушка: Complete при ID в диапазоне
возвращает \code{Task{ID:id,Title:"Отчет",Done:true},nil};
Delete возвращает nil. При ID вне диапазона методы заглушки запрещены.}
{ID в диапазоне: Complete возвращает заданный Task и nil,
Delete --- nil, соответствующий счетчик=1 и аргумент равен ID.
Вне диапазона: Complete дает Task{} и ErrID через errors.Is;
Delete дает ErrID через errors.Is; все счетчики=0. Проверка testing.}

\testcase{U04. Цель: сохранить порядок списка и представить отсутствие записей срезом длины 0 (F02)}
{Общий и краевой}
{\code{service.(*Service).List}}
{Вызов List без фильтра и параметров.}
{Три независимых возврата заглушки:
\code{[]Task{{ID:1,Title:"А",Done:false},{ID:2,Title:"Б",Done:true}},nil};
\code{[]Task{},nil}; \code{nil,nil}. БД отсутствует.}
{Первый набор: ровно две структуры в том же порядке, err=nil.
Второй и третий: \code{tasks!=nil}, \code{len(tasks)==0}, err=nil.
List заглушки вызван один раз. Поля и порядок сравниваются через testing.}

\testcase{U05. Цель: передать классификацию отказа Repository без данных Task (F05)}
{Негативный}
{Create, List, Complete, Delete сервиса}
{Create с «Отчет» (L=5); List без аргументов; Complete и Delete
с ID=9 (в диапазоне).}
{Независимые наборы: все четыре метода заглушки возвращают
\code{model.ErrStorage}; дополнительно Complete и Delete возвращают
\code{model.ErrNotFound}. Структура при отказе --- Task{}, список --- nil.}
{Для шести наборов errors.Is с заданной ошибкой равен true.
Create/Complete возвращают Task{}; List возвращает nil;
у Delete только error. Вызван ровно один соответствующий метод;
остальные счетчики=0. Сверка testing и errors.Is.}

\subsection{Интеграционные тесты}
\testcase{I01. Цель: запись и чтение через настоящий драйвер SQLite (F01, F02)}
{Общий и краевой}
{\code{repository.(*SQLiteRepository).Create}, \code{List}}
{Сначала List; затем Create с «Отчет» (L=5); затем List.
Длина названия в диапазоне 1--100.}
{F0; настоящий database/sql; заглушек нет.}
{Первый List: []Task{} и nil. Create:
\code{Task{ID:1,Title:"Отчет",Done:false}}, nil.
Второй List: срез из этого Task и nil.
Независимый \code{SELECT id,title,done FROM tasks ORDER BY id}
через database/sql возвращает ровно (1, «Отчет», 0);
\code{SELECT COUNT(*) FROM tasks} дает 1.
Строка title сравнивается побайтово.}

\testcase{I02. Цель: завершение повторяемо, удаление физическое (F03--F05)}
{Общий и негативный}
{Service + SQLiteRepository: Complete, Delete}
{Последовательно Complete(1), Complete(1), Delete(1), Delete(1).
ID=1 в диапазоне.}
{F1; настоящие сервис, репозиторий и БД, заглушек нет.}
{Оба Complete: \code{Task{ID:1,Title:"Отчет",Done:true}}, nil;
после каждого SELECT done WHERE id=1 дает 1 и COUNT(*)=1.
Первый Delete: nil, COUNT(*)=0. Второй: ErrNotFound через errors.Is,
COUNT(*) остается 0. SELECT и сравнения выполняет тест через database/sql.}

\testcase{I03. Цель: пройти HTTP-контракт жизненного цикла (F01--F04)}
{Общий}
{Полный стек; POST/GET/PATCH/DELETE из таблицы~\ref{tab:api}}
{POST \code{/api/tasks} с \code{{"title":"Отчет"}} (L=5);
GET списка; PATCH \code{/api/tasks/1/done}; DELETE \code{/api/tasks/1};
GET списка. Тела PATCH/DELETE пусты, ID=1.}
{F0; настоящий handler/service/repository;
запросы подаются через httptest.NewRecorder; Cookie нет.}
{Коды по порядку: 201, 200, 200, 204, 200.
Тела: Task(1,«Отчет»,false); массив из него;
Task(1,«Отчет»,true); 0 байт; [].
Проверка JSON и Content-Type по подразделу 1.3;
после DELETE независимый SELECT COUNT(*) дает 0.}

\testcase{I04. Цель: отказы входного HTTP не изменяют БД (F05)}
{Негативный и краевой}
{POST /api/tasks; PATCH /api/tasks/\{id\}/done;
маршрутизация PUT и неизвестного пути}
{Отдельные POST-тела: \code{{"title":""}},
\code{{"title":""}} с подстановкой T101 вместо пустой строки,
\code{{"title":}, {"title":null}, {"title":7}, {},
{"title":"А","extra":1}, {"title":"А"} {}}.
Дополнительно POST с телом \code{{"title":"А"}} и
\code{Content-Type: text/plain}.
PATCH с каждым ID: \code{0}, \code{-1}, \code{01}, \code{abc},
\code{2147483648}, \code{9}.
PUT /api/tasks без тела; GET /api/unknown.
Допуски: название L=1--100; ID и JSON по подразделам 1.2--1.3.}
{Для каждого набора отдельная F1; реальные слои;
кроме text/plain используется application/json.}
{POST с пустой строкой/T101: 400 и \code{{"error":"title_length"}}.
Шесть поврежденных/несоответствующих схеме JSON: 400 и
\code{{"error":"bad_json"}}; text/plain: 415 и \code{media_type}.
Пять ID вне синтаксиса/диапазона: 400 и \code{bad_id}; ID=9: 404 и \code{not_found}.
PUT: 405 и method, Allow=GET, POST; неизвестный путь: 404 и route.
Для каждого: JSON/Content-Type по 1.3, SQL возвращает только (1,«Отчет»,0),
COUNT(*)=1. Проверки httptest и database/sql.}

\testcase{I05. Цель: отказ БД преобразуется в HTTP 500 без служебных сведений (F05)}
{Негативный}
{Handler + Service + SQLiteRepository}
{Отдельно GET /api/tasks; POST /api/tasks с
\code{{"title":"Отчет"}} (L=5); PATCH /api/tasks/1/done;
DELETE /api/tasks/1. ID=1.}
{Для каждого набора создается F1, затем \code{db.Close()} возвращает nil;
handler продолжает использовать этот закрытый sql.DB.}
{Каждый ответ: 500, Content-Type по 1.3,
после декодирования ровно \code{{"error":"storage"}}.
Тело не содержит пути файла, текста SQL и сообщения \code{database is closed}.
httptest сравнивает код и структуру JSON.
Повторное открытие файла независимым sql.DB и SELECT подтверждают
единственную строку (1,«Отчет»,0).}

\subsection{Аттестационные тесты}
\testcase{A01. Цель: создать задачу и показать отказ для пустого ввода (F01, F05)}
{Общий и негативный}
{Форма GET /; \#title-input, \#create-button, \#error-banner}
{Открыть /, ввести «  Отчет  », нажать создание; затем ввести
три ASCII-пробела и повторить. После TrimSpace длины 5 и 0,
допуск L=1--100.}
{F0; настоящий сервер; HTML-форма не блокирует отправку пробелов
встроенной проверкой required. Cookie нет.}
{Первый POST: 201, Task(1,«Отчет»,false).
DOM: одна строка [data-id="1"], .title.textContent=«Отчет»,
.status.textContent=«Не выполнена», input.value="".
Второй POST: 400, \code{error=title_length};
\#error-banner.textContent=«Название: от 1 до 100 символов».
Число строк остается 1; SQL COUNT(*)=1 и title=«Отчет».
Playwright проверяет ответ, DOM в течение 2 с и SQL через тестовый драйвер.}

\testcase{A02. Цель: пустая страница и порядок записей совпадают с API (F02)}
{Краевой и общий}
{GET /; GET /api/tasks; \#tasks-body}
{На F0 открыть / и получить список API.
Затем через API создать \code{{"title":"Б"}},
\code{{"title":"А"}} (L=1 у обоих), перезагрузить страницу.}
{F0; отдельный настоящий сервер; обе записи создаются
последовательно, без других запросов записи.}
{До создания API: 200 и []; в \#tasks-body 0 строк.
Два POST: 201 и Task(1,«Б»,false), Task(2,«А»,false).
После перезагрузки API: 200 и массив этих двух задач в порядке ID 1,2;
DOM имеет 2 строки в том же порядке и тексты «Б», «А».
Проверка Playwright: массив data-id и textContent совпадает
с эталонными массивами ["1","2"] и ["Б","А"].}

\testcase{A03. Цель: завершенная задача сохраняется после остановки процесса (F03, F06)}
{Общий}
{.done-button; PATCH; процесс cmd/tasks; SQLite-файл}
{На странице F1 нажать завершение ID=1; затем независимым HTTP-клиентом
повторить PATCH /api/tasks/1/done без тела; остановить и запустить
сервер с тем же \code{DB_PATH}; снова открыть /.}
{F1; init и повторный старт по тайм-аутам подраздела 2.4;
файл не удаляется и не заменяется; заглушек нет.}
{Оба PATCH: 200 и Task(1,«Отчет»,true).
После первого ответа .status.textContent=«Выполнена»,
.done-button.disabled=true. После перезапуска GET списка: 200,
массив из той же задачи; DOM содержит одну строку с ID=1,
тем же названием, статусом и disabled=true.
SQL SELECT id,title,done дает ровно (1,«Отчет»,1).
Проверки Playwright, HTTP-клиент, database/sql.}

\testcase{A04. Цель: удаление отражается в UI, API и БД (F04, F05)}
{Общий и негативный}
{.delete-button; DELETE /api/tasks/1; GET списка}
{На F1 нажать удаление ID=1; затем выполнить повторный
DELETE /api/tasks/1 с пустым телом и GET /api/tasks. ID=1 в диапазоне.}
{F1; настоящий сервер и браузер, Cookie нет.}
{Первый DELETE: 204, 0 байт; DOM-строк с data-id="1" равно 0
в течение 2 с после ответа. Повтор: 404 и \code{{"error":"not_found"}}.
GET: 200 и []. SQL COUNT(*)=0.
Проверка Playwright/HTTP-клиентом и независимым SQL.}

\testcase{A05. Цель: показать пользователю отказ хранилища (F05)}
{Негативный}
{Форма создания, баннер, POST /api/tasks}
{На открытой странице ввести «Отчет» (L=5),
после подготовки отказа нажать создание.}
{F0. После загрузки страницы отдельным sql.DB выполнить
\code{DROP TABLE tasks}; это намеренная порча только тестового файла.
Сервер и браузер не перезапускаются; HTTP не подменяется.}
{POST: 500 и \code{{"error":"storage"}};
\#error-banner.textContent=«Хранилище недоступно»;
input.value=«Отчет», число строк в DOM=0.
Playwright проверяет ответ/DOM; независимый запрос
\code{SELECT COUNT(*) FROM sqlite_master WHERE type='table' AND name='tasks'}
дает 0. Отсутствие tasks не маскируется автоматическим созданием
схемы в обработчике запроса. После теста файл удаляется.}

\subsection{Специальные тесты}
\testcase{S01. Цель: выдержать 50 запросов чтения в секунду (N01)}
{Специальный, нагрузочный}
{GET /api/tasks; весь сервер}
{3000 плановых GET без тела за 60 с, 50 итераций/с,
1 запрос/итерацию; режим и VU по подразделу 2.5.1.}
{Пустая F0; отдельный опорный стенд; после прогрева файл вновь F0;
Cookie нет, иных клиентов и операций записи нет.}
{Каждый ответ: 200 и [], Content-Type по 1.3.
Сводка k6: N>=2940, N/60>=49 RPS, p95<=200 мс,
\code{dropped_iterations=0}, 5xx=0, \code{http_req_failed=0}, checks=1.
SQL COUNT(*)=0. Сверка счетчиков/thresholds и JSON-сводки k6,
независимый SELECT после прогона.}

\testcase{S02. Цель: выдержать 50 операций создания в секунду без потери записей (N01, F01)}
{Специальный, нагрузочный}
{POST /api/tasks; весь сервер и файловая SQLite}
{Каждый POST: \code{{"title":"Нагрузка"}} (L=8, допуск 1--100),
application/json. 3000 плановых запросов за 60 с, режим S01.}
{Новая F0 после прогрева; стенд как S01, Cookie нет.
Названия-дубликаты разрешены; иных писателей нет.}
{Каждый ответ: 201, поля ровно id/title/done,
целый ID в диапазоне, title=«Нагрузка», done=false.
Все пороги S01 по k6 соблюдены.
Пусть N201 --- число ответов 201: SQL COUNT(*) и COUNT(DISTINCT id)
равны N201; число строк с title='Нагрузка' AND done=0 также N201.
MIN(id)=1, MAX(id)=N201 при отсутствии неудавшихся INSERT.
Проверки k6 и независимыми SQL после завершения всех запросов.}

\testcase{S03. Цель: SQL-фрагмент остается значением title (N02)}
{Специальный, безопасность}
{POST /api/tasks; Repository.Create; GET списка}
{POST \code{{"title":"x'); DROP TABLE tasks;--"}}.
Название содержит 24 ASCII-символа, L=24 в диапазоне 1--100.}
{F1, настоящая SQLite; все SQL-значения передаются через параметры;
Cookie нет.}
{POST: 201, Task с ID=2, title точно равен переданной строке,
done=false. \code{SELECT COUNT(*) FROM tasks} дает 2;
\code{SELECT title FROM tasks WHERE id=2} совпадает со строкой побайтово;
строка ID=1 остается (1,«Отчет»,0). GET списка: 200 и эти две задачи
по ID ASC. Проверка HTTP-клиентом и database/sql.}

\testcase{S04. Цель: вставка в DOM через fetch не исполняет обработчик HTML (N03)}
{Специальный, безопасность}
{Форма создания; динамическое обновление .title}
{Ввести \code{<img src=x onerror=window.__xss=1>} (34 ASCII-символа,
L=34 в диапазоне), нажать создание.}
{F0; Playwright init-script задает \code{window.__xss=0} до загрузки;
учитывается число dialog-событий, изначально 0.}
{POST: 201, Task(1,переданный текст,false).
.title.textContent побайтово равен строке ввода;
\code{.title img} имеет count=0, \code{window.__xss=0},
число dialog-событий=0. Playwright проверяет после появления строки
и повторно спустя 500 мс наблюдения; SQL title WHERE id=1 совпадает
с исходной строкой. Указанный интервал ограничивает область наблюдения.}

\testcase{S05. Цель: начальный HTML экранирует сохраненный script (N03)}
{Специальный, безопасность}
{GET /; html/template; начальная .title}
{До навигации POST \code{{"title":"<script>window.__xss=1</script>"}}
(L=31, допуск 1--100), затем открыть / в новом контексте браузера.}
{F0; POST создает ID=1; init-script до навигации задает \code{window.__xss=0};
JavaScript страницы не должен заменять начальную таблицу до ее проверки.}
{POST: 201 и Task с точным названием, false.
GET /: 200 и HTML; HTML-парсер ответа видит в .title текст
\code{<script>window.__xss=1</script>}, дочерних script=0.
DOM после загрузки: та же строка, \code{.title script} count=0,
\code{window.__xss=0}, dialogs=0; повтор наблюдения через 500 мс.
Playwright и разбор HTML HTTP-ответа проверяют отдельно шаблон
и состояние браузера; SQL title WHERE id=1 равен входной строке.}

\subsection{Пример реализации блочного теста}
Ниже приведено содержимое \code{internal/service/service_test.go}
для модуля \code{example.com/tasks}. Фрагмент использует контракты
подраздела 2.1: \code{service.New}, методы Repository и ошибки model.
Он не требует testify или генерации моков. Приложение еще не реализовано;
после реализации контрактов файл запускается командой
\code{go test ./internal/service -run TestCreate -v}.
Русское слово «Отчет» записано Unicode-escape, поэтому листинг ASCII.
Листинг покрывает наборы U01/U02; U03--U05 реализуются отдельно
до перехода к интеграционному этапу.

\begin{lstlisting}
package service_test

import (
    "context"
    "errors"
    "strings"
    "testing"

    "example.com/tasks/internal/model"
    "example.com/tasks/internal/service"
)

type repoStub struct {
    createCalls int
    otherCalls int
    gotTitle string
}

func (r *repoStub) Create(_ context.Context,
    title string) (model.Task, error) {
    r.createCalls++
    r.gotTitle = title
    return model.Task{ID: 1, Title: title, Done: false}, nil
}
func (r *repoStub) List(context.Context) ([]model.Task, error) {
    r.otherCalls++
    return nil, errors.New("unexpected List")
}
func (r *repoStub) Complete(context.Context,
    int64) (model.Task, error) {
    r.otherCalls++
    return model.Task{}, errors.New("unexpected Complete")
}
func (r *repoStub) Delete(context.Context, int64) error {
    r.otherCalls++
    return errors.New("unexpected Delete")
}

func TestCreate(t *testing.T) {
    const title = "\u041e\u0442\u0447\u0435\u0442"
    cases := []struct {
        name, input, stored string
        wantErr error
        wantCalls int
    }{
        {"trim", "  " + title + "  ", title, nil, 1},
        {"lower", "\u042f", "\u042f", nil, 1},
        {"upper", strings.Repeat("A", 100),
            strings.Repeat("A", 100), nil, 1},
        {"empty", "", "", model.ErrTitle, 0},
        {"spaces", "   ", "", model.ErrTitle, 0},
        {"too_long", strings.Repeat("A", 101),
            "", model.ErrTitle, 0},
        {"bad_utf8", string([]byte{0xff}),
            "", model.ErrTitle, 0},
    }
    for _, tc := range cases {
        t.Run(tc.name, func(t *testing.T) {
            repo := &repoStub{}
            svc := service.New(repo)
            got, err := svc.Create(context.Background(), tc.input)
            if !errors.Is(err, tc.wantErr) {
                t.Fatalf("error: got %v, want %v", err, tc.wantErr)
            }
            want := model.Task{}
            if tc.wantErr == nil {
                want = model.Task{ID: 1, Title: tc.stored}
            }
            if got != want {
                t.Fatalf("task: got %#v, want %#v", got, want)
            }
            if repo.createCalls != tc.wantCalls ||
                repo.otherCalls != 0 || repo.gotTitle != tc.stored {
                t.Fatalf("repo state: %#v", repo)
            }
        })
    }
}
\end{lstlisting}

\subsection{Покрытие и плановые значения}\label{sec:coverage}
Покрытие операторов измеряется только для service, repository и handler:
\begin{lstlisting}[language={}]
go test -count=1 -covermode=atomic \
  -coverpkg=./internal/service,./internal/repository,./internal/handler \
  -coverprofile=coverage.out ./internal/...
go tool cover -func=coverage.out
go tool cover -html=coverage.out -o coverage.html
\end{lstlisting}

Пусть $s_i$ --- число операторов в инструментированном блоке профиля,
$c_i$ --- счетчик исполнения этого блока. Тогда
\[
 C_{\mathrm{stmt}}=
 \frac{\sum_i s_i\,\mathbf{1}[c_i>0]}{\sum_i s_i}\cdot100\%.
\]
План: $C_{\mathrm{stmt}}\geq85\%$ для трех пакетов в совокупности,
отдельно service не ниже 90\%. model без исполняемой логики,
cmd со сборкой зависимостей, тестовые файлы и внешние зависимости
в знаменатель не входят. Их исключение явно сохраняется в отчете.

\code{go test -coverprofile} не вычисляет покрытие ветвей.
Для него после реализации составляется реестр с позицией в исходнике,
условием, исходом true/false либо case/default, ID подтеста и отметкой
исполнения. Выполнение исхода подтверждается трассировкой отладчика
или временным счетчиком в диагностической сборке; счетчики не входят
в поставляемую программу. Покрытие операторов двух блоков само по себе
не доказывает покрытие всех составных логических условий.
\[
 C_{\mathrm{branch}}=
 \frac{B_{\mathrm{executed}}}{B_{\mathrm{total}}}\cdot100\%.
\]
План: 100\% исходов проверок названия/ID и обработки результатов
Repository в service; не ниже 80\% исходов во всех трех пакетах.
Это покрытие решений; MC/DC и покрытие всех путей не заявляются.
При отсутствии измерения ветвей записывается «Не измерено»,
а не число из coverage.out. Итоговые числа вносятся только в раздел 6.

Покрытие требований рассчитывается отдельно от кода:
\[
 C_{\mathrm{req}}=\frac{R_{\mathrm{linked}}}{R_{\mathrm{total}}}\cdot100\%,
 \qquad
 C_{\mathrm{req,pass}}=\frac{R_{\mathrm{passed}}}{R_{\mathrm{total}}}\cdot100\%.
\]
$R_{\mathrm{linked}}$ --- число требований с тестами, содержащими
проверки всех их условий; $R_{\mathrm{passed}}$ --- число требований,
для которых все связанные обязательные наборы имеют PASS.
Знаменатель фиксирован: 6 функциональных и 3 нефункциональных,
то есть $R_{\mathrm{total}}=9$. Плановое покрытие спецификациями ---
9/9=100\%; это оценка полноты плана, а не результат выполнения.
Для аттестации отдельно: 6/6=100\% по F01--F06.

\begin{longtable}{|P{20mm}|P{63mm}|P{69mm}|}
\caption{Матрица трассировки требований}\label{tab:trace}\\\hline
Требование & U и I & A или S \\\hline\endfirsthead
\multicolumn{3}{l}{Продолжение таблицы \thetable}\\\hline
Требование & U и I & A или S \\\hline\endhead
F01 & U01, U02, I01, I03 & A01 \\\hline
F02 & U04, I01, I03 & A02 \\\hline
F03 & U03, I02, I03 & A03 \\\hline
F04 & U03, I02, I03 & A04 \\\hline
F05 & U02, U03, U05, I02, I04, I05 & A01, A04, A05 \\\hline
F06 & --- & A03 \\\hline
N01 & --- & S01, S02 \\\hline
N02 & --- & S03 \\\hline
N03 & --- & S04, S05 \\\hline
\end{longtable}

Если измерение не достигает порога, добавляются тесты с ID,
продолжающими последовательность своего уровня, и обновляется матрица.
Порог не понижается задним числом. Ни 20 спецификаций, ни 100\% покрытия
требований не являются доказательством отсутствия дефектов.

\mainsection{Журнал тестирования (шаблон)}
Заполняется после проведения тестирования.
Одна строка относится к тесту или именованному параметризованному набору;
повторные прогоны учитываются в количестве запусков. ID дефекта
имеет вид BUG-001. Даты записываются с временем и часовым поясом
в связанной записи прогона. Фактические статусы и артефакты сохраняются
в отчете инструмента с тем же ID прогона.

\begin{longtable}{|P{21mm}|P{23mm}|P{29mm}|P{24mm}|P{20mm}|P{29mm}|}
\caption{Журнал тестирования}\\\hline
Дата & Тестировщик & Объект тестирования & ID выполненного теста &
Количество запусков & Перечень найденных ошибок (ID ошибок) \\\hline\endfirsthead
\multicolumn{6}{l}{Продолжение таблицы \thetable}\\\hline
Дата & Тестировщик & Объект тестирования & ID выполненного теста &
Количество запусков & Перечень найденных ошибок (ID ошибок) \\\hline\endhead
\rule{0pt}{16mm} & & & & & \\\hline
\rule{0pt}{16mm} & & & & & \\\hline
\rule{0pt}{16mm} & & & & & \\\hline
\rule{0pt}{16mm} & & & & & \\\hline
\rule{0pt}{16mm} & & & & & \\\hline
\end{longtable}

\mainsection{Журнал найденных ошибок (шаблон)}
Заполняется после проведения тестирования.
В связанном описании BUG указываются commit, версия плана,
среда, серьезность, статус, шаги воспроизведения и артефакты.
В таблице приводятся наблюдаемые значения: код HTTP, JSON,
поля Task, строки SQL или DOM, а не общая оценка поведения.

\begin{longtable}{|P{25mm}|P{23mm}|P{20mm}|P{40mm}|P{40mm}|}
\caption{Журнал найденных ошибок}\\\hline
Номер отчета об ошибке & Дата составления & Номер (ID) теста &
Ожидаемый результат & Фактический результат \\\hline\endfirsthead
\multicolumn{5}{l}{Продолжение таблицы \thetable}\\\hline
Номер отчета об ошибке & Дата составления & Номер (ID) теста &
Ожидаемый результат & Фактический результат \\\hline\endhead
\rule{0pt}{20mm} & & & & \\\hline
\rule{0pt}{20mm} & & & & \\\hline
\rule{0pt}{20mm} & & & & \\\hline
\rule{0pt}{20mm} & & & & \\\hline
\end{longtable}

\mainsection{Результаты (шаблон)}
Заполняется после проведения тестирования.
\subsection{Идентификация исследованной версии}
\begin{longtable}{|P{68mm}|P{84mm}|}
\caption{Паспорт итогового прогона}\\\hline
Поле & Значение \\\hline
Commit приложения, версия плана & \rule{0pt}{10mm} \\\hline
Дата, время, часовой пояс, исполнитель & \rule{0pt}{10mm} \\\hline
ОС, CPU, RAM, носитель БД & \rule{0pt}{10mm} \\\hline
Go, драйвер, SQLite, CGO, компилятор C & \rule{0pt}{10mm} \\\hline
Node.js, Playwright, Chromium, k6 & \rule{0pt}{10mm} \\\hline
Пути к логам, trace, coverage, сводке k6 & \rule{0pt}{10mm} \\\hline
\end{longtable}

\subsection{Оценка результатов и критерии приемки}
Единица учета --- обязательный набор входных данных; у U03 это
каждая пара метод/ID. Состав наборов фиксируется до запуска.
Доля пройденных наборов:
\[
 P=\frac{N_{\mathrm{PASS}}}
 {N_{\mathrm{PASS}}+N_{\mathrm{FAIL}}+N_{\mathrm{BLOCKED}}+
 N_{\mathrm{NOT\ RUN}}}\cdot100\%.
\]
Для итоговой приемки требуется $P=100\%$; BLOCKED и NOT RUN
не исключаются из знаменателя. Сводка по ID содержит PASS только
при PASS всех наборов соответствующей спецификации.

\begin{longtable}{|P{40mm}|P{80mm}|P{30mm}|}
\caption{Показатели приемки: значения заполняются после тестирования}\\\hline
Показатель & Плановый критерий & Измерено \\\hline\endfirsthead
\multicolumn{3}{l}{Продолжение таблицы \thetable}\\\hline
Показатель & Плановый критерий & Измерено \\\hline\endhead
U / I / A / S & Все обязательные наборы каждого уровня PASS;
P=100\%, FAIL=0, BLOCKED=0, NOT RUN=0 & \rule{0pt}{12mm} \\\hline
Операторы & Совокупно не менее 85\%; service не менее 90\% & \\\hline
Ветви & Проверки ввода и отказы Repository в service: 100\%;
три пакета: не менее 80\% & \\\hline
Требования & Creq,pass=100\% для 9 требований;
аттестация F01--F06: 100\% & \\\hline
Гонки & go test -race: код 0, DATA RACE=0 & \\\hline
Нагрузка, каждый прогон & RPS не менее 49; N не менее 2940;
p95 не более 200 мс; 5xx=0; dropped=0; failed=0; checks=1 & \\\hline
Безопасность & S03: обе записи сохранены; S04/S05:
счетчик XSS=0, внедренных DOM-узлов=0, dialogs=0 & \\\hline
Открытые дефекты & Блокирующих, критических и существенных 0;
незначительные перечислены и согласованы проверяющим & \\\hline
\end{longtable}

\subsection{Оценка качества и решение о приемке}
Заполняется после проведения тестирования.
\begin{longtable}{|P{68mm}|P{84mm}|}
\caption{Итоговая оценка}\\\hline
Позиция & Заключение по фактическим данным \\\hline
Реализованные требования и исключения & \rule{0pt}{14mm} \\\hline
Открытые дефекты и влияние на работу & \rule{0pt}{14mm} \\\hline
Ограничения стенда и полноты проверок & \rule{0pt}{14mm} \\\hline
Решение о приемке и основание & \rule{0pt}{14mm} \\\hline
Проверяющий, подпись, дата & \rule{0pt}{14mm} \\\hline
\end{longtable}

\struct{ЗАКЛЮЧЕНИЕ}
План задает поведение будущего приложения и способы наблюдения
результата на уровне возвращаемых значений, HTTP, SQLite и DOM.
Определены порядок интеграции, условия переходов, пороги нагрузки,
проверка защиты двух точек отображения названия и измерение покрытия.
Журналы и итоговые измерения не заполнены: реализация приложения
и выполнение тестов относятся к следующему этапу работы.

\struct{СПИСОК ИСПОЛЬЗОВАННЫХ ИСТОЧНИКОВ}
\begin{enumerate}[label=\arabic*.]
\item ГОСТ 7.32-2001. Система стандартов по информации, библиотечному
и издательскому делу. Отчет о научно-исследовательской работе.
Структура и правила оформления. С изменением №~1.
\item The Go Programming Language. The cover story [Электронный ресурс].
URL: \url{https://go.dev/blog/cover} (дата обращения: 15.11.2026).
\item SQLite. SQL Language Expressions: Parameters [Электронный ресурс].
URL: \url{https://www.sqlite.org/lang_expr.html}
(дата обращения: 15.11.2026).
\item Grafana k6. Thresholds [Электронный ресурс].
URL: \url{https://grafana.com/docs/k6/latest/using-k6/thresholds/}
(дата обращения: 15.11.2026).
\end{enumerate}
\end{document}
